% !TeX root = ../初等数学笔记.tex
% 由主文档通过 \input 载入；请编译主文档。
\setcounter{chapter}{15} % 保留原书章号。
\chapter{立体几何与度量}
\label{chap:em-solid}
空间中的线面关系比平面复杂，尤其要区分相交与投影相交。解决度量问题的基本策略，是先确定空间位置，再截取包含目标长度或角的恰当平面，使用已经建立的平面几何。

\section{空间对象与相互位置}
本章采用欧氏空间的基本关联事实：不共线三点确定一个平面；直线上两点在某平面内，则整条直线在该平面内；两个不同平面若相交，其公共部分为一条直线。两不同直线可相交、平行或不共面，最后一种称为\emph{异面直线}（skew lines）。

直线与平面的关系为包含、相交一点或平行且无交点；两不同平面相交或平行。透视图中两线相交只表示投影相交，不能推出原空间直线有公共点。例如立方体一条底边与不在同一侧面上的竖棱可能为异面直线。

\section{平行与垂直的判定}
常用欧氏空间基本定理如下：平面外的一条直线若平行于平面内某一直线，则与该平面平行；一个平面内两条相交直线分别平行于另一平面内两条相交直线，则两平面平行。不能把“相交直线”减为两条平行直线，因为它们只提供一个方向。

直线垂直于平面，意为它垂直于经过垂足的每一条平面内直线；判定时只需检查其中两条相交直线。两平面垂直的常用充分条件，是其中一平面包含另一平面的一条垂线。这些判定作为欧氏空间的基础定理使用，完整的空间公理化证明不在本书展开；使用时必须确认有关直线确实属于指定平面。

例如长方体一个顶点出发的三条棱两两垂直，任意一条棱垂直于另外两棱所确定的面。已知只垂直于平面内一条直线却不够，因为该直线在平面内也有垂线。

\section{空间角、距离与截面}
异面直线夹角通过平移方向到共同点定义，通常取不超过 $90^\circ$ 的角。非垂直直线与平面的夹角是它与平面上正射影的夹角；垂直情形定义为 $90^\circ$。二面角的大小则在垂直于棱的截面中测量两半平面的夹角，范围为 $[0,180^\circ]$。

点到平面的距离为垂线段长度。若 $H$ 为垂足、$Q$ 为平面上其他点，则 $PQ^2=PH^2+HQ^2\ge PH^2$，故垂线段最短。长方体边长 $a,b,c>0$ 时，先求底面对角线 $\sqrt{a^2+b^2}$，再在含体对角线的矩形截面应用勾股定理，得体对角线 $\sqrt{a^2+b^2+c^2}$。两次所用直角必须来自线面关系，而非示意图外观。

\section{棱柱、棱锥与旋转体}
棱柱具有两个平行全等底面，侧棱平行等长；棱锥由一个底面多边形和一个不在其平面内的顶点组成。圆柱、圆锥可由矩形、直角三角形绕恰当边旋转形成。高是底面之间或顶点到底面所在平面的垂直距离，母线和斜高一般不是高。

直圆柱半径 $r$、高 $h$ 的侧面展开为边长 $2\pi r,h$ 的矩形，侧面积为 $2\pi rh$；直圆锥底半径 $r$、母线长 $\ell$ 的侧面展开为半径 $\ell$、弧长 $2\pi r$ 的扇形，侧面积为 $\pi r\ell$。球面是到球心距离为 $r$ 的点集，球体还包含内部。长度、面积、体积在比例 $k>0$ 的相似缩放下分别乘以 $k,k^2,k^3$。

\section{体积公式与球的度量}
采用通常体积的可加性、相似缩放规律与截面比较原理：两个立体夹在相同高度的平行平面间，若每个高度的截面积相等，则体积相等。该原理需要体积对逼近的相容性，本书作为度量基础使用，不把有限几张截面图当作证明。

底面积为 $B$、高为 $h$ 的柱体体积为 $Bh$。先注意同高同底面积的锥体在相同相对高度处截面等面积，所以截面比较保证它们等体积。把三棱柱 $ABC$--$A^{\prime}B^{\prime}C^{\prime}$ 分为四面体 $ABCC^{\prime}$、$ABB^{\prime}C^{\prime}$、$AA^{\prime}B^{\prime}C^{\prime}$：前两者共享底面 $ABC^{\prime}$，顶点 $C,B^{\prime}$ 关于该底面具有相等距离，因为侧面平行四边形的两对角线互相平分；后两者对公共底面 $AB^{\prime}C^{\prime}$ 同理等高。因此三者等体积，各为柱体的三分之一，得到锥体体积 $Bh/3$。一般底面按三角形分割后由可加性推广。平行于底面截得的锥台，两底面积为 $B_1,B_2$、高为 $h$，由相似锥体体积相减得
\[
 V=\frac h3(B_1+B_2+\sqrt{B_1B_2}).
\]

半径 $r$ 的半球在离球心平面高度 $z\in[0,r]$ 处截面积为 $\pi(r^2-z^2)$。半径、高都为 $r$ 的圆柱挖去顶点在底面中心、顶面为底的圆锥，其同高度截面积也为 $\pi r^2-\pi z^2$。因此半球体积为 $\pi r^3-\pi r^3/3=2\pi r^3/3$，整球体积为 $4\pi r^3/3$。

球面面积为 $4\pi r^2$。这里引用球面带面积定理：高度为 $h$ 的球面带面积为 $2\pi rh$；其严格证明需用内接圆台侧面积的极限定义曲面面积，不由体积公式在纯代数意义下自动推出。取整球高度 $2r$ 即得结果，相关逼近细节可在已有的圆与球度量专题中继续学习。

\section{例题、习题与学习检查}
\begin{example}
直圆锥的底半径为 $3$、高为 $4$，求母线、表面积和体积。
\end{example}
\begin{solve}
轴截面的一半是直角三角形，所以母线为 $5$。侧面积为 $15\pi$，加上底面积 $9\pi$，表面积为 $24\pi$；体积为 $\pi\cdot3^2\cdot4/3=12\pi$。面积使用平方单位，体积使用立方单位，两类结果不能混比。
\end{solve}
\begin{enumerate}
\item 长方体边长为 $2,3,6$，求体对角线。答：$7$。
\item 两相似球半径比为 $2:3$，表面积与体积比分别是多少？答：$4:9$ 与 $8:27$。
\item 高为 $H$ 的锥体在离顶点 $H/2$ 处作平行底面的截面，顶部小锥体占总体积多少？答：$1/8$。
\end{enumerate}

